% !TEX root = ../mainfile.tex
% Table 6.1: the dollar's response by event and horizon. % !TEX root = ../mainfile.tex
% Table 6.1: the dollar's response by event and horizon. % !TEX root = ../mainfile.tex
% Table 6.1: the dollar's response by event and horizon. % !TEX root = ../mainfile.tex
% Table 6.1: the dollar's response by event and horizon. \input{content/tab_es_dollar_horizon} at the end of the section that discusses it.
% Split 2026-09-09 from tab_event_study_results.tex (hand-maintained; numbers generated by
% Code/03_event_study.py and Code/04_event_study_w50.py, constant-mean model, estimation
% window [-140,-21]; [-170,-51] for h=50; verified against Data/processed/event_study/*.csv).

\begin{table}[htbp]
\centering
\begin{threeparttable}
\caption{The U.S.\ dollar's response by event and horizon}
\label{tab:es_dollar_horizon}
\small
\setlength{\tabcolsep}{4.5pt}
\begin{tabular*}{\textwidth}{@{\extracolsep{\fill}} l rrr rrr @{}}
\toprule
 & \multicolumn{3}{c}{Headline windows} & \multicolumn{3}{c}{Longer horizons} \\
\cmidrule(lr){2-4}\cmidrule(lr){5-7}
Event & \multicolumn{1}{c}{$(-1,+1)$} & \multicolumn{1}{c}{$(0,+1)$} & \multicolumn{1}{c}{$(0,+5)$} & \multicolumn{1}{c}{$(0,+10)$} & \multicolumn{1}{c}{$(0,+20)$} & \multicolumn{1}{c@{}}{$(0,+50)$} \\
\midrule
\multicolumn{7}{l}{\emph{Panel A. U.S.-originated trade-policy shock}}\\
\quad Liberation Day        & $-1.70^{***}$ & $-1.44^{***}$ & $-0.10\phantom{^{***}}$       & $-3.00^{***}$ & $-3.98^{***}$ & $-6.34^{***}$ \\
\quad\quad Tariff pause     & $-1.48^{**}\phantom{^{*}}$  & $-1.48^{***}$ & $-3.07^{***}$ & $-3.44^{***}$ & $-4.51^{***}$ & $-6.74^{***}$ \\
\addlinespace
\multicolumn{7}{l}{\emph{Panel B. U.S.-involved military/oil-supply shocks}}\\
\quad Twelve-day war        & $-0.28\phantom{^{***}}$       & $+0.04\phantom{^{***}}$       & $+0.88\phantom{^{***}}$       & $-0.13\phantom{^{***}}$       & $+0.35\phantom{^{***}}$       & $-0.12\phantom{^{***}}$       \\
\quad Hormuz crisis         & $+1.34^{***}$ & $+1.39^{***}$ & $+1.53^{***}$ & $+2.10^{***}$ & $+3.25^{***}$ & $-0.13\phantom{^{***}}$       \\
\quad\quad Hormuz ceasefire & $-1.33^{***}$ & $-1.14^{***}$ & $-1.72^{***}$ & $-1.65^{**}\phantom{^{*}}$  & $-1.47\phantom{^{***}}$       & $-0.04\phantom{^{***}}$       \\
\addlinespace
\multicolumn{7}{l}{\emph{Panel C. External control}}\\
\quad Ukraine invasion      & $+0.45\phantom{^{***}}$       & $+0.49\phantom{^{***}}$       & $+1.03^{*}\phantom{^{**}}$   & $+1.65^{**}\phantom{^{*}}$  & $+0.84\phantom{^{***}}$       & $+3.12^{**}\phantom{^{*}}$  \\
\bottomrule
\end{tabular*}
\begin{tablenotes}[flushleft]\footnotesize
\item \textit{Notes.} Cumulative abnormal returns of the broad nominal U.S.\ dollar index, in percent, constant-mean model. Estimation windows are as in Section~\ref{sec:method_event}, with the fifty-day column estimated on $[-170,-51]$. The three headline windows carry the conclusions. The longer horizons trace how long a response lasts (Section~\ref{sec:es_robust}). Positive values denote dollar appreciation. The tariff pause and the Hormuz ceasefire are sub-events within their parent windows. The ceasefire's estimation window ends on 9~March and contains the first six trading days of the crisis, so its cells are read with corresponding caution. The May~2026 U.S.\ strikes and the June~2026 memorandum are not analysed (Section~\ref{sec:es_robust}). $^{*}$, $^{**}$, $^{***}$ denote significance at the 10\%, 5\% and 1\% level.
% superseded 2026-09-09, kept for reference:
% \item \textit{Notes.} Cumulative abnormal returns of the broad nominal U.S.\ dollar
% index (FRED \texttt{DTWEXBGS}), in percent, from a constant-mean model. Estimation
% window $[-140,-21]$ trading days for $h\le 20$, ending where the event window begins;
% $[-170,-51]$ for $h{=}50$ (pushed back so it does not overlap the wider event window).
% Inference rests on the headline windows; the longer horizons gauge durability and
% persistence but absorb confounding events and are read as robustness
% (Section~\ref{sec:method_event}). Positive values denote dollar appreciation. Tariff
% pause and Hormuz ceasefire are sub-events within their parent windows; ``Hormuz ceasefire'' is the two-week ceasefire of 7~April 2026, distinct from the memorandum of understanding of 17~June 2026. Two further events are not reported. The estimation window of the 25~May 2026 U.S.\ strikes contains the Hormuz crisis, so their abnormal returns would be measured against a benchmark the crisis itself inflated. The Hormuz ceasefire shares this overlap in a milder form, with 26 of its 120 estimation days inside the crisis window, and is reported nonetheless as a within-crisis sub-event, its cells read with corresponding caution. The memorandum falls too close to the sample end of 30~June 2026 for a full event window.
% $^{*}$, $^{**}$, $^{***}$ denote significance at the 10\%, 5\% and 1\% level.
\end{tablenotes}
\end{threeparttable}
\end{table}
 at the end of the section that discusses it.
% Split 2026-09-09 from tab_event_study_results.tex (hand-maintained; numbers generated by
% Code/03_event_study.py and Code/04_event_study_w50.py, constant-mean model, estimation
% window [-140,-21]; [-170,-51] for h=50; verified against Data/processed/event_study/*.csv).

\begin{table}[htbp]
\centering
\begin{threeparttable}
\caption{The U.S.\ dollar's response by event and horizon}
\label{tab:es_dollar_horizon}
\small
\setlength{\tabcolsep}{4.5pt}
\begin{tabular*}{\textwidth}{@{\extracolsep{\fill}} l rrr rrr @{}}
\toprule
 & \multicolumn{3}{c}{Headline windows} & \multicolumn{3}{c}{Longer horizons} \\
\cmidrule(lr){2-4}\cmidrule(lr){5-7}
Event & \multicolumn{1}{c}{$(-1,+1)$} & \multicolumn{1}{c}{$(0,+1)$} & \multicolumn{1}{c}{$(0,+5)$} & \multicolumn{1}{c}{$(0,+10)$} & \multicolumn{1}{c}{$(0,+20)$} & \multicolumn{1}{c@{}}{$(0,+50)$} \\
\midrule
\multicolumn{7}{l}{\emph{Panel A. U.S.-originated trade-policy shock}}\\
\quad Liberation Day        & $-1.70^{***}$ & $-1.44^{***}$ & $-0.10\phantom{^{***}}$       & $-3.00^{***}$ & $-3.98^{***}$ & $-6.34^{***}$ \\
\quad\quad Tariff pause     & $-1.48^{**}\phantom{^{*}}$  & $-1.48^{***}$ & $-3.07^{***}$ & $-3.44^{***}$ & $-4.51^{***}$ & $-6.74^{***}$ \\
\addlinespace
\multicolumn{7}{l}{\emph{Panel B. U.S.-involved military/oil-supply shocks}}\\
\quad Twelve-day war        & $-0.28\phantom{^{***}}$       & $+0.04\phantom{^{***}}$       & $+0.88\phantom{^{***}}$       & $-0.13\phantom{^{***}}$       & $+0.35\phantom{^{***}}$       & $-0.12\phantom{^{***}}$       \\
\quad Hormuz crisis         & $+1.34^{***}$ & $+1.39^{***}$ & $+1.53^{***}$ & $+2.10^{***}$ & $+3.25^{***}$ & $-0.13\phantom{^{***}}$       \\
\quad\quad Hormuz ceasefire & $-1.33^{***}$ & $-1.14^{***}$ & $-1.72^{***}$ & $-1.65^{**}\phantom{^{*}}$  & $-1.47\phantom{^{***}}$       & $-0.04\phantom{^{***}}$       \\
\addlinespace
\multicolumn{7}{l}{\emph{Panel C. External control}}\\
\quad Ukraine invasion      & $+0.45\phantom{^{***}}$       & $+0.49\phantom{^{***}}$       & $+1.03^{*}\phantom{^{**}}$   & $+1.65^{**}\phantom{^{*}}$  & $+0.84\phantom{^{***}}$       & $+3.12^{**}\phantom{^{*}}$  \\
\bottomrule
\end{tabular*}
\begin{tablenotes}[flushleft]\footnotesize
\item \textit{Notes.} Cumulative abnormal returns of the broad nominal U.S.\ dollar index, in percent, constant-mean model. Estimation windows are as in Section~\ref{sec:method_event}, with the fifty-day column estimated on $[-170,-51]$. The three headline windows carry the conclusions. The longer horizons trace how long a response lasts (Section~\ref{sec:es_robust}). Positive values denote dollar appreciation. The tariff pause and the Hormuz ceasefire are sub-events within their parent windows. The ceasefire's estimation window ends on 9~March and contains the first six trading days of the crisis, so its cells are read with corresponding caution. The May~2026 U.S.\ strikes and the June~2026 memorandum are not analysed (Section~\ref{sec:es_robust}). $^{*}$, $^{**}$, $^{***}$ denote significance at the 10\%, 5\% and 1\% level.
% superseded 2026-09-09, kept for reference:
% \item \textit{Notes.} Cumulative abnormal returns of the broad nominal U.S.\ dollar
% index (FRED \texttt{DTWEXBGS}), in percent, from a constant-mean model. Estimation
% window $[-140,-21]$ trading days for $h\le 20$, ending where the event window begins;
% $[-170,-51]$ for $h{=}50$ (pushed back so it does not overlap the wider event window).
% Inference rests on the headline windows; the longer horizons gauge durability and
% persistence but absorb confounding events and are read as robustness
% (Section~\ref{sec:method_event}). Positive values denote dollar appreciation. Tariff
% pause and Hormuz ceasefire are sub-events within their parent windows; ``Hormuz ceasefire'' is the two-week ceasefire of 7~April 2026, distinct from the memorandum of understanding of 17~June 2026. Two further events are not reported. The estimation window of the 25~May 2026 U.S.\ strikes contains the Hormuz crisis, so their abnormal returns would be measured against a benchmark the crisis itself inflated. The Hormuz ceasefire shares this overlap in a milder form, with 26 of its 120 estimation days inside the crisis window, and is reported nonetheless as a within-crisis sub-event, its cells read with corresponding caution. The memorandum falls too close to the sample end of 30~June 2026 for a full event window.
% $^{*}$, $^{**}$, $^{***}$ denote significance at the 10\%, 5\% and 1\% level.
\end{tablenotes}
\end{threeparttable}
\end{table}
 at the end of the section that discusses it.
% Split 2026-09-09 from tab_event_study_results.tex (hand-maintained; numbers generated by
% Code/03_event_study.py and Code/04_event_study_w50.py, constant-mean model, estimation
% window [-140,-21]; [-170,-51] for h=50; verified against Data/processed/event_study/*.csv).

\begin{table}[htbp]
\centering
\begin{threeparttable}
\caption{The U.S.\ dollar's response by event and horizon}
\label{tab:es_dollar_horizon}
\small
\setlength{\tabcolsep}{4.5pt}
\begin{tabular*}{\textwidth}{@{\extracolsep{\fill}} l rrr rrr @{}}
\toprule
 & \multicolumn{3}{c}{Headline windows} & \multicolumn{3}{c}{Longer horizons} \\
\cmidrule(lr){2-4}\cmidrule(lr){5-7}
Event & \multicolumn{1}{c}{$(-1,+1)$} & \multicolumn{1}{c}{$(0,+1)$} & \multicolumn{1}{c}{$(0,+5)$} & \multicolumn{1}{c}{$(0,+10)$} & \multicolumn{1}{c}{$(0,+20)$} & \multicolumn{1}{c@{}}{$(0,+50)$} \\
\midrule
\multicolumn{7}{l}{\emph{Panel A. U.S.-originated trade-policy shock}}\\
\quad Liberation Day        & $-1.70^{***}$ & $-1.44^{***}$ & $-0.10\phantom{^{***}}$       & $-3.00^{***}$ & $-3.98^{***}$ & $-6.34^{***}$ \\
\quad\quad Tariff pause     & $-1.48^{**}\phantom{^{*}}$  & $-1.48^{***}$ & $-3.07^{***}$ & $-3.44^{***}$ & $-4.51^{***}$ & $-6.74^{***}$ \\
\addlinespace
\multicolumn{7}{l}{\emph{Panel B. U.S.-involved military/oil-supply shocks}}\\
\quad Twelve-day war        & $-0.28\phantom{^{***}}$       & $+0.04\phantom{^{***}}$       & $+0.88\phantom{^{***}}$       & $-0.13\phantom{^{***}}$       & $+0.35\phantom{^{***}}$       & $-0.12\phantom{^{***}}$       \\
\quad Hormuz crisis         & $+1.34^{***}$ & $+1.39^{***}$ & $+1.53^{***}$ & $+2.10^{***}$ & $+3.25^{***}$ & $-0.13\phantom{^{***}}$       \\
\quad\quad Hormuz ceasefire & $-1.33^{***}$ & $-1.14^{***}$ & $-1.72^{***}$ & $-1.65^{**}\phantom{^{*}}$  & $-1.47\phantom{^{***}}$       & $-0.04\phantom{^{***}}$       \\
\addlinespace
\multicolumn{7}{l}{\emph{Panel C. External control}}\\
\quad Ukraine invasion      & $+0.45\phantom{^{***}}$       & $+0.49\phantom{^{***}}$       & $+1.03^{*}\phantom{^{**}}$   & $+1.65^{**}\phantom{^{*}}$  & $+0.84\phantom{^{***}}$       & $+3.12^{**}\phantom{^{*}}$  \\
\bottomrule
\end{tabular*}
\begin{tablenotes}[flushleft]\footnotesize
\item \textit{Notes.} Cumulative abnormal returns of the broad nominal U.S.\ dollar index, in percent, constant-mean model. Estimation windows are as in Section~\ref{sec:method_event}, with the fifty-day column estimated on $[-170,-51]$. The three headline windows carry the conclusions. The longer horizons trace how long a response lasts (Section~\ref{sec:es_robust}). Positive values denote dollar appreciation. The tariff pause and the Hormuz ceasefire are sub-events within their parent windows. The ceasefire's estimation window ends on 9~March and contains the first six trading days of the crisis, so its cells are read with corresponding caution. The May~2026 U.S.\ strikes and the June~2026 memorandum are not analysed (Section~\ref{sec:es_robust}). $^{*}$, $^{**}$, $^{***}$ denote significance at the 10\%, 5\% and 1\% level.
% superseded 2026-09-09, kept for reference:
% \item \textit{Notes.} Cumulative abnormal returns of the broad nominal U.S.\ dollar
% index (FRED \texttt{DTWEXBGS}), in percent, from a constant-mean model. Estimation
% window $[-140,-21]$ trading days for $h\le 20$, ending where the event window begins;
% $[-170,-51]$ for $h{=}50$ (pushed back so it does not overlap the wider event window).
% Inference rests on the headline windows; the longer horizons gauge durability and
% persistence but absorb confounding events and are read as robustness
% (Section~\ref{sec:method_event}). Positive values denote dollar appreciation. Tariff
% pause and Hormuz ceasefire are sub-events within their parent windows; ``Hormuz ceasefire'' is the two-week ceasefire of 7~April 2026, distinct from the memorandum of understanding of 17~June 2026. Two further events are not reported. The estimation window of the 25~May 2026 U.S.\ strikes contains the Hormuz crisis, so their abnormal returns would be measured against a benchmark the crisis itself inflated. The Hormuz ceasefire shares this overlap in a milder form, with 26 of its 120 estimation days inside the crisis window, and is reported nonetheless as a within-crisis sub-event, its cells read with corresponding caution. The memorandum falls too close to the sample end of 30~June 2026 for a full event window.
% $^{*}$, $^{**}$, $^{***}$ denote significance at the 10\%, 5\% and 1\% level.
\end{tablenotes}
\end{threeparttable}
\end{table}
 at the end of the section that discusses it.
% Split 2026-09-09 from tab_event_study_results.tex (hand-maintained; numbers generated by
% Code/03_event_study.py and Code/04_event_study_w50.py, constant-mean model, estimation
% window [-140,-21]; [-170,-51] for h=50; verified against Data/processed/event_study/*.csv).

\begin{table}[htbp]
\centering
\begin{threeparttable}
\caption{The U.S.\ dollar's response by event and horizon}
\label{tab:es_dollar_horizon}
\small
\setlength{\tabcolsep}{4.5pt}
\begin{tabular*}{\textwidth}{@{\extracolsep{\fill}} l rrr rrr @{}}
\toprule
 & \multicolumn{3}{c}{Headline windows} & \multicolumn{3}{c}{Longer horizons} \\
\cmidrule(lr){2-4}\cmidrule(lr){5-7}
Event & \multicolumn{1}{c}{$(-1,+1)$} & \multicolumn{1}{c}{$(0,+1)$} & \multicolumn{1}{c}{$(0,+5)$} & \multicolumn{1}{c}{$(0,+10)$} & \multicolumn{1}{c}{$(0,+20)$} & \multicolumn{1}{c@{}}{$(0,+50)$} \\
\midrule
\multicolumn{7}{l}{\emph{Panel A. U.S.-originated trade-policy shock}}\\
\quad Liberation Day        & $-1.70^{***}$ & $-1.44^{***}$ & $-0.10\phantom{^{***}}$       & $-3.00^{***}$ & $-3.98^{***}$ & $-6.34^{***}$ \\
\quad\quad Tariff pause     & $-1.48^{**}\phantom{^{*}}$  & $-1.48^{***}$ & $-3.07^{***}$ & $-3.44^{***}$ & $-4.51^{***}$ & $-6.74^{***}$ \\
\addlinespace
\multicolumn{7}{l}{\emph{Panel B. U.S.-involved military/oil-supply shocks}}\\
\quad Twelve-day war        & $-0.28\phantom{^{***}}$       & $+0.04\phantom{^{***}}$       & $+0.88\phantom{^{***}}$       & $-0.13\phantom{^{***}}$       & $+0.35\phantom{^{***}}$       & $-0.12\phantom{^{***}}$       \\
\quad Hormuz crisis         & $+1.34^{***}$ & $+1.39^{***}$ & $+1.53^{***}$ & $+2.10^{***}$ & $+3.25^{***}$ & $-0.13\phantom{^{***}}$       \\
\quad\quad Hormuz ceasefire & $-1.33^{***}$ & $-1.14^{***}$ & $-1.72^{***}$ & $-1.65^{**}\phantom{^{*}}$  & $-1.47\phantom{^{***}}$       & $-0.04\phantom{^{***}}$       \\
\addlinespace
\multicolumn{7}{l}{\emph{Panel C. External control}}\\
\quad Ukraine invasion      & $+0.45\phantom{^{***}}$       & $+0.49\phantom{^{***}}$       & $+1.03^{*}\phantom{^{**}}$   & $+1.65^{**}\phantom{^{*}}$  & $+0.84\phantom{^{***}}$       & $+3.12^{**}\phantom{^{*}}$  \\
\bottomrule
\end{tabular*}
\begin{tablenotes}[flushleft]\footnotesize
\item \textit{Notes.} Cumulative abnormal returns of the broad nominal U.S.\ dollar index, in percent, constant-mean model. Estimation windows are as in Section~\ref{sec:method_event}, with the fifty-day column estimated on $[-170,-51]$. The three headline windows carry the conclusions. The longer horizons trace how long a response lasts (Section~\ref{sec:es_robust}). Positive values denote dollar appreciation. The tariff pause and the Hormuz ceasefire are sub-events within their parent windows. The ceasefire's estimation window ends on 9~March and contains the first six trading days of the crisis, so its cells are read with corresponding caution. The May~2026 U.S.\ strikes and the June~2026 memorandum are not analysed (Section~\ref{sec:es_robust}). $^{*}$, $^{**}$, $^{***}$ denote significance at the 10\%, 5\% and 1\% level.
% superseded 2026-09-09, kept for reference:
% \item \textit{Notes.} Cumulative abnormal returns of the broad nominal U.S.\ dollar
% index (FRED \texttt{DTWEXBGS}), in percent, from a constant-mean model. Estimation
% window $[-140,-21]$ trading days for $h\le 20$, ending where the event window begins;
% $[-170,-51]$ for $h{=}50$ (pushed back so it does not overlap the wider event window).
% Inference rests on the headline windows; the longer horizons gauge durability and
% persistence but absorb confounding events and are read as robustness
% (Section~\ref{sec:method_event}). Positive values denote dollar appreciation. Tariff
% pause and Hormuz ceasefire are sub-events within their parent windows; ``Hormuz ceasefire'' is the two-week ceasefire of 7~April 2026, distinct from the memorandum of understanding of 17~June 2026. Two further events are not reported. The estimation window of the 25~May 2026 U.S.\ strikes contains the Hormuz crisis, so their abnormal returns would be measured against a benchmark the crisis itself inflated. The Hormuz ceasefire shares this overlap in a milder form, with 26 of its 120 estimation days inside the crisis window, and is reported nonetheless as a within-crisis sub-event, its cells read with corresponding caution. The memorandum falls too close to the sample end of 30~June 2026 for a full event window.
% $^{*}$, $^{**}$, $^{***}$ denote significance at the 10\%, 5\% and 1\% level.
\end{tablenotes}
\end{threeparttable}
\end{table}
